
\documentclass[11pt]{article}

\usepackage[margin=1in]{geometry}
\usepackage{booktabs}
\usepackage{array}
\usepackage{multirow}
\usepackage{graphicx}
\usepackage{amsmath}
\usepackage{amssymb}
\usepackage{natbib}
\usepackage{microtype}
\usepackage{hyperref}
\usepackage{xcolor}
\usepackage{enumitem}
\usepackage{caption}

\hypersetup{
  colorlinks=true,
  linkcolor=blue,
  citecolor=blue,
  urlcolor=blue
}

\title{When Does History Start Matter?\\
Model-Specific Sensitivity and Global Robustness in Retail Forecast Evaluation}

\author{Jinyong Ha\\
Graduate School of Policy Studies, Korea University\\
Seoul, Republic of Korea\\
\texttt{hjy075@korea.ac.kr} \quad \texttt{hjy075@gmail.com}}
\date{August 2026}

\emergencystretch=2em
\begin{document}
\maketitle

\begin{abstract}
Retail forecasting pipelines must decide where a product--store history begins. A common convention is to remove leading zero sales and start a series at its first positive observation, but such zeros may reflect genuine active periods rather than pre-availability. We study how this history-start convention changes forecast evaluation while holding the target observation fixed. Using Visuelle 2.0 and FreshRetailNet-50K, we compare a lifecycle-aligned or stock-grounded history with a first-positive-trimmed history across nine canonical statistical forecasting methods. In Visuelle, the affected subset is rare (262 of 106,850 series; 0.245\%), yet seven of nine methods exhibit non-zero bootstrap effects and local rankings are unstable (mean Kendall's $\tau=0.483$), while the top method changes in only 10.6\% of bootstrap samples. In FreshRetailNet-50K, 903 of 50,000 series (1.806\%) contain stock-confirmed active-zero observations before first sale. Effects are smaller and heterogeneous: HistoricAverage has a positive 95\% product-cluster bootstrap interval excluding zero, optimized SES has a negative interval excluding zero, and point-effect signs vary non-monotonically across evaluation origins; the corrected winner-change probability is 33.4\%. In both datasets, benchmark-level scores and rankings remain largely stable, consistent with the limited prevalence of affected histories. We also show that history-dependent MASE/RMSSE denominators can reverse the apparent direction of an effect. The results motivate separating local forecast sensitivity, evaluation eligibility, benchmark propagation, and metric-normalization sensitivity when assessing preprocessing conventions.
\end{abstract}

\section{Introduction}

Forecast evaluation is usually discussed in terms of models, horizons, metrics, and test-set construction. A more basic design choice often receives less attention: \emph{where does the historical series begin?} This question matters in retail because product--store histories can begin with zeros for several different reasons. A zero may represent a product that was not yet available, a product that was available but sold nothing, or a data-collection convention. Treating all leading zeros as equivalent can therefore change the information supplied to a forecasting method.

First-positive initialization is not merely hypothetical. Oracle Supply Chain Planning documentation states that its demand forecasting process removes leading zero demand before the first positive demand point \citep{oracle2026leadingzeros}. Related academic work also preprocesses intermittent-demand series by removing initial zeros and taking the first non-zero demand as the initial observation when the meaning of initial zeros is ambiguous \citep{li2023feature}. Such choices can be operationally reasonable, but they are also \emph{outcome-conditioned}: the history start is determined by the location of the first positive realization.

This paper asks a narrow empirical question: \emph{when first-positive initialization differs from a lifecycle-aligned or stock-grounded history start, how much does the history-start convention alter forecast performance and method rankings, and do those local changes propagate to benchmark-level conclusions?} The focus is evaluation validity rather than proposing a new forecasting model.

We treat the history-start convention as an explicit component of forecast-evaluation design rather than as a neutral preprocessing detail. To isolate its consequences, we separate three evaluation channels. First, the \textbf{local common-support effect} compares forecasts for the same series, target time, and realized target while changing only the history start. Second, the \textbf{eligibility effect} measures whether trimming removes early forecast origins under a minimum-history requirement before forecast accuracy is evaluated. Third, the \textbf{global propagation effect} asks whether local changes are sufficiently prevalent to alter benchmark-level scores, rankings, or winners. For history-dependent scaled metrics, we additionally distinguish \textbf{metric-normalization sensitivity} by separating changes in the forecasts from changes in the scaling denominator.

Two public datasets provide complementary identification. Visuelle 2.0 contains short lifecycle histories for fast-fashion products \citep{skenderi2022visuelle}; its market-delivery setting gives a natural lifecycle-aligned baseline, but affected series are rare. FreshRetailNet-50K contains 50,000 store--product series with explicit stockout annotations \citep{wang2025freshretail}; it does not document product-launch timing at the observation-window start, but it allows us to identify leading zero-sales days on which the item was not recorded as stocked out. We therefore use FreshRetailNet-50K as a \emph{mechanism replication}, not as a product-launch replication.

Our main findings are:
\begin{enumerate}[leftmargin=1.5em]
    \item History-start sensitivity is \textbf{model-specific}. In Visuelle, seven of nine canonical methods have product-cluster bootstrap confidence intervals for $\Delta$MAE that exclude zero. In FreshRetailNet-50K, only HistoricAverage and optimized SES remain individually distinguishable from zero after the estimand-corrected bootstrap, with opposite signs.
    \item Local \textbf{ranking sensitivity} survives even when the deterministic winner is stable. Corrected product-cluster bootstraps give mean Kendall's $\tau$ of 0.483 in Visuelle and 0.813 in FreshRetailNet-50K, with winner changes in 10.6\% and 33.4\% of replicates, respectively.
    \item FreshRetailNet-50K reveals strong \textbf{temporal heterogeneity}. The sign of the trimming effect changes repeatedly across evaluation origins rather than decaying monotonically with more history.
    \item First-positive initialization creates a distinct \textbf{coverage effect}: 15.2\% of eligible treated origins are lost in Visuelle and 1.46\% in the strict FreshRetailNet-50K sample under the primary minimum-history rule.
    \item Local instability does \textbf{not} imply global benchmark failure. Because exposure is rare, benchmark-level score changes are very small and global rankings are nearly unchanged.
    \item Scaled metrics introduce a separate measurement channel. When MASE/RMSSE denominators are recomputed from each convention-specific history, the apparent direction of an effect can differ from the fixed-denominator result.
\end{enumerate}

Taken together, these results characterize history-start sensitivity rather than supporting a universal claim that first-positive trimming is harmful. The central empirical finding is that model-specific local sensitivity, rank changes, and eligibility loss can coexist with globally robust benchmark conclusions when affected histories are rare. A further measurement result is that history-dependent scaling can change, and in some cases reverse, the apparent direction of the preprocessing effect. These distinctions make history-start semantics, exposure prevalence, evaluation eligibility, benchmark propagation, and metric normalization explicit components of forecast-evaluation design.

\section{Related Work}

\subsection{Intermittent demand forecasting}

Intermittent demand contains many zero observations and irregular positive arrivals, motivating methods that treat demand sizes and inter-demand intervals separately. Croston's method is the classical starting point \citep{croston1972}. Subsequent work identified bias in Croston-style estimates and proposed corrections, including the Syntetos--Boylan approximation \citep{syntetos2005accuracy}. The TSB method updates demand probability every period, which allows forecasts to react to extended zero runs and obsolescence \citep{teunter2011tsb}. Temporal aggregation has also been used to reduce intermittence, including ADIDA \citep{nikolopoulos2011adida} and multi-level forecast combinations \citep{petropoulos2015combinations}.

Our study does not compare these methods to establish a new accuracy ranking. Instead, we use a heterogeneous panel of canonical implementations to test whether the \emph{same method} changes when its input history begins at a different point. This distinction matters because the history-start convention can interact with each method's initialization and state updates.

\subsection{Leading zeros and outcome-conditioned history starts}

Leading-zero removal appears in both forecasting systems and empirical research. Current Oracle demand-plan documentation describes removal of leading zero demand before the first positive point \citep{oracle2026leadingzeros}. In an intermittent-demand forecasting study using RAF and M5 data, \citet{li2023feature} explicitly remove initial zeros because the data do not reveal whether those zeros mean zero demand or that a product was not yet in stores.

This ambiguity motivates our identification strategy. Rather than assuming that all leading zeros are either ``inactive'' or ``true zero demand,'' we exploit lifecycle information in Visuelle and stock-status information in FreshRetailNet-50K. The resulting comparison is deliberately conditional: it evaluates series for which first-positive trimming removes observations that the alternative convention would retain.

\subsection{Benchmark validity and forecast metrics}

Retail forecasting competitions have emphasized the importance of evaluation design and representativeness. The M5 competition formalized large-scale retail forecast evaluation across hierarchical sales series \citep{makridakis2022m5}; subsequent work examined how representative those series are of other retail datasets \citep{theodorou2022representativeness}. Our question is complementary: even within a fixed benchmark, can a preprocessing convention change which observations are admitted into history and thereby alter model comparisons?

Forecast metrics add another layer. \citet{hyndman2006mase} proposed MASE to enable scale-free comparisons using an in-sample naive-error denominator. That denominator is itself history-dependent. When the research intervention changes the history start, recomputing the denominator under each convention mixes two effects: a change in forecasts and a change in scale normalization. We therefore report both fixed-denominator and convention-specific scaled metrics.

\section{Data}

\subsection{Visuelle 2.0}

Visuelle 2.0 is a fast-fashion benchmark containing 5,355 products across six seasons and item--shop sales histories \citep{skenderi2022visuelle}. We obtained the dataset through the official Visuelle 2.0 project page, which provides access via a short form \citep{visuelle_dataset_page}. The released \texttt{sales.csv} table used here contains 106,850 product--store series from 110 stores, each observed over a 12-week window. We exclude series with negative weekly sales from the primary global analysis, leaving 103,492 series. We use only the sales table; images, customer-level purchase data, and other modalities are not used in this study.

We define a strict affected series as one for which (i) the first positive sale occurs after week 0, (ii) every preceding weekly sale is exactly zero, and (iii) the 12-week series contains no negative sales. This yields 262 series (0.245\% of all series), spanning 245 products and 87 stores. The baseline history begins at week 0 of the provided lifecycle window; the alternative history begins at the first positive weekly sale.

Visuelle provides the cleaner lifecycle-oriented comparison but low prevalence. It therefore offers strong local identification while making large global benchmark effects unlikely by construction.

\subsection{FreshRetailNet-50K}

FreshRetailNet-50K contains 50,000 store--product series with detailed sales and stockout annotations from fresh retail \citep{wang2025freshretail}. We use the official Dingdong-Inc release on Hugging Face \citep{freshretail_dataset_card}, specifically the public training split with 4.5 million daily rows and exactly 90 days per series. The dataset paper reports 898 stores in 18 major cities. Our analysis uses only \texttt{store\_id}, \texttt{product\_id}, \texttt{dt}, \texttt{sale\_amount}, and \texttt{stock\_hour6\_22\_cnt}.

The start of the 90-day window is \emph{not} documented as a product launch date. Accordingly, we do not call the baseline ``release-aligned.'' Instead, we define a strict stock-grounded sample using the daily stockout variable. A series is treated when (i) its first positive sale occurs after day 0, (ii) all preceding sales are exactly zero, (iii) all preceding days have \texttt{stock\_hour6\_22\_cnt == 0}, indicating no recorded stockout hours during 06:00--22:00, and (iv) the series has no negative sales.

Of 50,000 series, 3,256 (6.512\%) have a delayed first positive sale, but only 903 (1.806\%) satisfy the strict stock-grounded definition. These 903 series span 229 products and 563 stores. Thus only about 27.7\% of delayed-first-positive series have fully stock-confirmed active-zero leading periods under our definition. This is an important descriptive result: neither ``all leading zeros are pre-availability'' nor ``all leading zeros are genuine active zeros'' is supported.

\begin{table}[t]
\centering
\caption{Datasets and affected samples. FreshRetailNet-50K is a stock-grounded mechanism replication rather than a launch-date replication.}
\label{tab:data}
\small
\begin{tabular}{lrrrr}
\toprule
Dataset & All series & Strict affected & Share & Unique products \\
\midrule
Visuelle 2.0 & 106,850 & 262 & 0.245\% & 245 \\
FreshRetailNet-50K & 50,000 & 903 & 1.806\% & 229 \\
\bottomrule
\end{tabular}
\end{table}

\section{Experimental Design}

\subsection{Paired common-support estimand}

Let $y_{i,t}$ denote the realized demand for series $i$ at target time $t$. Let $s_i$ be the grounded history start and $f_i$ the first-positive time. For forecasting method $m$, the two one-step forecasts are
\begin{align}
\hat y^{G}_{i,t,m} &= \mathcal{F}_m(y_{i,s_i:t-1}),\\
\hat y^{T}_{i,t,m} &= \mathcal{F}_m(y_{i,f_i:t-1}),
\end{align}
where $G$ denotes grounded/lifecycle-aligned history and $T$ denotes first-positive trimming. The target $y_{i,t}$ is identical in both conditions.

For absolute error, we first average the paired error difference within each series over its common-support target set $\mathcal{T}_i$, and then average equally across the $N$ affected series:
\begin{equation}
\Delta_m =
\frac{1}{N}
\sum_{i=1}^{N}
\left[
\frac{1}{|\mathcal{T}_i|}
\sum_{t\in\mathcal{T}_i}
\left(
|y_{i,t}-\hat y^T_{i,t,m}|
-
|y_{i,t}-\hat y^G_{i,t,m}|
\right)
\right].
\end{equation}
This is a series-weighted estimand: each product--store series receives equal weight regardless of how many common-support origins it contributes. Positive $\Delta_m$ means trimming worsens MAE; negative values mean trimming improves MAE.

The primary minimum-history requirement is three observations in both conditions. Visuelle additionally checks minimum histories of 2 and 4. FreshRetailNet-50K checks 2, 3, 4, and 7 observations. FreshRetailNet-50K uses prespecified target days
\[
\{4,7,14,21,28,42,56,70,84,89\},
\]
which allows the history-start effect to be examined from very short to long histories.

\subsection{Forecasting methods}

We use nine canonical implementations from StatsForecast 2.1.1 \citep{garza2022statsforecast}:
Naive, HistoricAverage, optimized simple exponential smoothing (SESOpt), CrostonClassic, CrostonOptimized, CrostonSBA, TSB, ADIDA, and IMAPA. TSB uses $\alpha_d=\alpha_p=0.2$ in the primary analysis; Visuelle additionally checks 0.1 and 0.3. Forecasts are truncated below at zero where applicable.

These methods span last-observation, mean-based, exponential-smoothing, Croston-style, probability-update, and temporal-aggregation families. The purpose is not to identify a state-of-the-art winner, but to test whether sensitivity generalizes across forecasting mechanisms.

\subsection{Uncertainty and rank stability}

For each method we first average forecast errors within each store--product series across eligible origins. We then use a product-cluster bootstrap with 5,000 replicates. Products are sampled with replacement as clusters, but all store-series belonging to sampled products are pooled, preserving the same \emph{series-weighted} estimand as the reported point estimate.

For each replicate we compute $\Delta_m$, model ranks under both history conventions, Kendall's $\tau$, and the winning method. This corrected procedure is important: a preliminary bootstrap that averaged products equally targeted a different product-weighted estimand. The final results reported here use cluster resampling with series weighting throughout.

\subsection{Eligibility and global propagation}

The common-support experiment deliberately conditions on targets forecastable under both histories. Trimming can also delay the earliest forecast origin. For a minimum history $h$, we define
\begin{equation}
\text{Coverage loss}
=
1-\frac{N_T(h)}{N_G(h)},
\end{equation}
where $N_G$ and $N_T$ are the numbers of eligible forecast origins under the grounded and trimmed histories.

Finally, we propagate the history convention to the full benchmark. Unaffected series have identical forecasts under both conditions; affected series use their alternative histories. Visuelle global evaluation uses target weeks 8--11. FreshRetailNet-50K uses target days 28, 56, and 89. We report global MAE, WAPE, and deterministic rank changes. FreshRetailNet-50K also includes a sensitivity analysis that trims all 3,256 delayed-first-positive series regardless of stock confirmation; this ``blanket'' scenario is operational rather than causally grounded.

\subsection{Scaled metrics}

MASE scales absolute forecast error by an in-sample naive-error denominator \citep{hyndman2006mase}. We report two variants:
\begin{itemize}[leftmargin=1.5em]
    \item \textbf{Fixed denominator}: both forecast conditions use the grounded-history scale. This isolates forecast changes.
    \item \textbf{Convention-specific denominator}: each condition computes its scale from its own history. This captures forecast changes plus metric-normalization changes.
\end{itemize}
We apply the same distinction to RMSSE. This separation prevents a history-start intervention from silently changing both the prediction and the measurement scale.

\section{Results}

\subsection{Visuelle: strong local sensitivity, stable winner}

Figure~\ref{fig:local} summarizes the corrected product-cluster bootstrap effects. In Visuelle, seven of nine methods have 95\% intervals that exclude zero, all in the direction of higher MAE after trimming. The largest changes occur for CrostonClassic ($\Delta$MAE $=0.165$, 95\% CI $[0.113,0.218]$), CrostonSBA ($0.135$, $[0.087,0.184]$), and CrostonOptimized ($0.086$, $[0.059,0.114]$). TSB, HistoricAverage, IMAPA, and ADIDA also have 95\% intervals entirely above zero. SESOpt has a small, uncertain effect and Naive is invariant by construction because the last observation is unchanged.

Although IMAPA remains the deterministic top method under both conditions, the full ranking changes materially. TSB moves from rank 3 to 5, CrostonSBA from 5 to 8, CrostonClassic from 6 to 9, SESOpt from 8 to 3, and Naive from 9 to 6. The corrected cluster bootstrap yields mean Kendall's $\tau=0.483$ (95\% interval $[0.222,0.722]$); 99.5\% of replicates have $\tau<0.8$. The winning method changes in 10.6\% of replicates. Across minimum-history settings 2, 3, and 4, the deterministic winner remains IMAPA while Kendall's $\tau$ ranges from 0.278 to 0.389, confirming that winner stability does not imply rank stability.

\begin{figure}[t]
\centering
\begin{minipage}{0.49\linewidth}
\centering
\includegraphics[width=\linewidth]{figures/visuelle_local_effects.png}
\end{minipage}
\hfill
\begin{minipage}{0.49\linewidth}
\centering
\includegraphics[width=\linewidth]{figures/freshretail_local_effects.png}
\end{minipage}
\caption{Local common-support effects under the corrected series-weighted product-cluster bootstrap. Positive $\Delta$MAE means first-positive trimming worsens accuracy. Visuelle effects are broadly positive; FreshRetailNet-50K effects are smaller and heterogeneous.}
\label{fig:local}
\end{figure}

\begin{table*}[t]
\centering
\caption{Corrected local effects and deterministic rank changes. Confidence intervals come from 5,000 product-cluster bootstrap replicates targeting the series-weighted estimand.}
\label{tab:local}
\scriptsize
\begin{tabular}{lccc ccc}
\toprule
& \multicolumn{3}{c}{Visuelle 2.0} & \multicolumn{3}{c}{FreshRetailNet-50K} \\
\cmidrule(lr){2-4}\cmidrule(lr){5-7}
Model & $\Delta$MAE & 95\% CI & Rank $G\rightarrow T$ & $\Delta$MAE & 95\% CI & Rank $G\rightarrow T$\\
\midrule
IMAPA & 0.034644 & [0.018029, 0.051943] & 1$\rightarrow$1 & 0.000925 & [-0.000647, 0.002508] & 3$\rightarrow$4 \\
ADIDA & 0.032805 & [0.015261, 0.050300] & 2$\rightarrow$2 & 0.001049 & [-0.000562, 0.002744] & 2$\rightarrow$3 \\
TSB & 0.057222 & [0.030294, 0.084530] & 3$\rightarrow$5 & 0.000147 & [-0.001890, 0.002182] & 1$\rightarrow$1 \\
HistoricAverage & 0.039108 & [0.023306, 0.054923] & 4$\rightarrow$4 & 0.002538 & [0.000871, 0.004601] & 7$\rightarrow$8 \\
CrostonSBA & 0.135055 & [0.086725, 0.184388] & 5$\rightarrow$8 & -0.000013 & [-0.002716, 0.002959] & 4$\rightarrow$2 \\
CrostonClassic & 0.165110 & [0.113383, 0.217732] & 6$\rightarrow$9 & 0.002374 & [-0.000506, 0.005578] & 5$\rightarrow$7 \\
CrostonOptimized & 0.086098 & [0.059272, 0.113902] & 7$\rightarrow$7 & 0.001808 & [-0.000005, 0.003660] & 6$\rightarrow$5 \\
SESOpt & 0.003845 & [-0.017766, 0.027473] & 8$\rightarrow$3 & -0.002115 & [-0.003846, -0.000452] & 8$\rightarrow$6 \\
Naive & 0 & [0,0] & 9$\rightarrow$6 & 0 & [0,0] & 9$\rightarrow$9 \\
\bottomrule
\end{tabular}
\end{table*}

\subsection{FreshRetailNet-50K: smaller effects and temporal heterogeneity}

FreshRetailNet-50K provides a different picture. Point effects are much smaller, and only two methods have corrected 95\% intervals excluding zero. HistoricAverage worsens ($\Delta$MAE $=0.002538$, 95\% CI $[0.000871,0.004601]$), while SESOpt \emph{improves} ($-0.002115$, $[-0.003846,-0.000452]$). This directly rejects the stronger hypothesis that first-positive trimming is uniformly harmful.

The deterministic winner, TSB, remains first under both conditions, but the order below the winner changes: CrostonSBA moves from rank 4 to 2, CrostonClassic from 5 to 7, and SESOpt from 8 to 6. Under the corrected cluster bootstrap, mean Kendall's $\tau$ is 0.813 (95\% interval $[0.611,0.944]$), and 43.2\% of replicates have $\tau<0.8$. The winner changes in 33.4\% of replicates. Release/grounded winner probabilities are diffuse---TSB 35.8\%, ADIDA 34.8\%, CrostonSBA 19.7\%, and IMAPA 9.1\%---and shift under trimming, with TSB rising to 46.7\%.

A distinctive FreshRetailNet-50K result is descriptive temporal heterogeneity in the point effects. Figure~\ref{fig:time} shows $\Delta$MAE over the ten prespecified target days. At day 4, most non-naive methods worsen substantially: for example CrostonClassic rises by 16.8\%, CrostonOptimized by 21.1\%, HistoricAverage by 17.8\%, and TSB by 11.8\%. At day 7, the direction reverses for the same families: CrostonClassic improves by 11.7\%, CrostonSBA by 13.1\%, and TSB by 9.5\%. Across models, the point-effect profile is non-monotonic, with two to six sign changes over the target grid. Thus the observed effect does not simply decay as history length grows; it interacts with the evaluation origin and model state. These origin-specific profiles are descriptive: we do not attach separate bootstrap confidence intervals to each target day.

\begin{figure}[t]
\centering
\includegraphics[width=0.88\linewidth]{figures/freshretail_target_day_profile.png}
\caption{FreshRetailNet-50K $\Delta$MAE across prespecified evaluation origins. The point-effect profile is non-monotonic and often changes sign, especially in short histories; origin-specific uncertainty is not shown.}
\label{fig:time}
\end{figure}

\subsection{Eligibility effects}

Table~\ref{tab:coverage} separates common-support accuracy from forecast eligibility. In Visuelle, first-positive initialization removes 359 forecast origins under the primary minimum-history rule: 15.22\% of treated origins but only 0.0385\% of all benchmark origins. The treated loss increases from 13.70\% at minimum history 2 to 17.13\% at minimum history 4.

In strict FreshRetailNet-50K, 1,146 treated origins are lost, corresponding to 1.46\% of treated origins and 0.0263\% globally. The lower treated loss is consistent with the empirical delay distribution: at least 75\% of strict series have only a one-day gap before first positive sale. In the broader blanket sensitivity scenario, trimming all 3,256 delayed-first-positive series produces a treated coverage loss of 2.80\% and a global loss of 0.1825\%.

\begin{table}[t]
\centering
\caption{Eligibility loss under the primary minimum-history rule ($h=3$).}
\label{tab:coverage}
\small
\begin{tabular}{lrrr}
\toprule
Dataset / scenario & Lost origins & Treated loss & Global loss \\
\midrule
Visuelle strict & 359 & 15.22\% & 0.0385\% \\
FreshRetail strict & 1,146 & 1.46\% & 0.0263\% \\
FreshRetail blanket sensitivity & 7,938 & 2.80\% & 0.1825\% \\
\bottomrule
\end{tabular}
\end{table}

\subsection{Global propagation is small}

Local sensitivity does not materially invalidate either full benchmark. For Visuelle, global evaluation over weeks 8--11 preserves the entire nine-model ranking. The largest MAE change is only 0.000551 for CrostonClassic (1.342678 to 1.343229).

For the strict FreshRetailNet-50K scenario over days 28, 56, and 89, global changes are also tiny. The largest absolute MAE change is about $4.47\times10^{-5}$ (CrostonClassic). Only CrostonSBA and CrostonOptimized exchange ranks 5 and 6. Under the broader blanket sensitivity, the largest absolute MAE change is approximately 0.000799, while the deterministic rank ordering remains unchanged.

These results follow naturally from prevalence. A preprocessing convention can have a meaningful effect conditional on being exposed to it, yet have little influence on a benchmark average when affected histories are rare. Consequently, claims about local evaluation validity and claims about benchmark-level invalidity must be separated.

\begin{figure}[t]
\centering
\includegraphics[width=0.70\linewidth]{figures/local_rank_stability.png}
\caption{Corrected local rank stability. Kendall's $\tau$ summarizes agreement between model rankings; winner-change probability measures whether the best model changes within the same product-cluster bootstrap replicate.}
\label{fig:rank}
\end{figure}

\subsection{Metric normalization can change the conclusion}

History-dependent scaled metrics create an additional source of sensitivity. In Visuelle, SESOpt has a fixed-denominator MASE change of $-0.0155$, suggesting improvement after trimming, but the convention-specific MASE change is $+0.0930$. In FreshRetailNet-50K the same method changes from $-0.0124$ under the fixed denominator to $+0.0328$ under convention-specific scaling. FreshRetail CrostonClassic similarly changes from $-0.0093$ to $+0.0354$.

The sign reversal is not a forecasting contradiction. The two calculations answer different questions. Fixed scaling asks whether the forecasts changed relative to a common benchmark scale. Convention-specific scaling allows the preprocessing intervention to alter both the forecast and the denominator. When the purpose is to study preprocessing sensitivity, we recommend reporting the fixed-denominator result as primary and the convention-specific result as a separate measurement-sensitivity analysis.

\section{Discussion}

\subsection{History-start sensitivity is not a universal bias}

The original motivating intuition could have supported a broad claim that first-positive trimming biases retail forecasting benchmarks. The data do not support that claim. Visuelle shows strong and mostly harmful local effects, whereas FreshRetailNet-50K shows smaller, mixed-sign effects and repeated sign reversals over evaluation time. A more defensible interpretation is \emph{history-start sensitivity}: the convention changes the information set, and different methods react differently.

This distinction is important conceptually. First-positive trimming is not automatically wrong. If leading zeros truly represent pre-availability periods, removing them may be appropriate. Conversely, if the product was active and available, those zeros are valid observations about demand occurrence. The correct history start is therefore a data-semantic question, not merely a forecasting hyperparameter.

\subsection{Winner stability and rank stability are different}

A second lesson is that evaluating only the winning model can hide instability. In Visuelle the deterministic winner remains IMAPA across all tested minimum-history settings even though the rest of the ranking is substantially reordered. In FreshRetailNet-50K, TSB remains the deterministic winner, yet one-third of corrected bootstrap replicates switch winners between history conventions. Benchmark reports should therefore distinguish at least three concepts: score sensitivity, rank sensitivity, and winner sensitivity.

\subsection{Local and global conclusions should be reported separately}

The two datasets converge most clearly on the local/global distinction. Conditional effects can be large enough to matter for affected product--store histories while remaining too rare to move benchmark-wide averages. This resembles a general representativeness issue in benchmark interpretation \citep{theodorou2022representativeness}: aggregate benchmark conclusions depend not only on the magnitude of an effect but also on how prevalent the relevant series type is.

For benchmark authors, we recommend reporting the prevalence of preprocessing-sensitive histories alongside accuracy results. For practitioners, the local subset may be more relevant than the benchmark average if newly launched, sparse, or recently activated items constitute an operationally important segment.

\subsection{Practical recommendations}

Our results suggest five simple practices for retail forecast evaluation:
\begin{enumerate}[leftmargin=1.5em]
    \item \textbf{Declare the history-start convention.} State whether leading zeros are retained, removed, or interpreted using availability metadata.
    \item \textbf{Use external state information when available.} Release dates, assortment status, and stockout indicators help distinguish pre-availability zeros from active-zero demand.
    \item \textbf{Separate common-support accuracy from eligibility.} A preprocessing rule can improve reported accuracy partly by making difficult early origins ineligible.
    \item \textbf{Report local and global effects.} A large affected-subset effect does not establish benchmark-wide instability.
    \item \textbf{Fix scaled-metric denominators for sensitivity studies.} If the intervention changes history, recomputing MASE/RMSSE scales can conflate forecast sensitivity with normalization sensitivity.
\end{enumerate}

\section{Limitations}

This study has several limitations. First, only two public retail datasets are examined. Visuelle provides a lifecycle-oriented history but has very low affected prevalence; FreshRetailNet-50K has stronger stock-status evidence but its observation-window start is not a documented product launch date. Second, the experiment focuses on one-step statistical forecasting methods rather than modern global neural forecasters. This is intentional for identification and computational reproducibility, but future work could test whether learned global models are more or less sensitive to history-start conventions. Third, FreshRetailNet-50K results are evaluated on a prespecified grid of target days rather than every possible origin, and the origin-specific profiles are descriptive because separate confidence intervals are not estimated for each target day. Fourth, the bootstrap clusters by product to preserve dependence among stores carrying the same product, but it does not explicitly model residual dependence among different products within the same store. Fifth, stock confirmation uses the released \texttt{stock\_hour6\_22\_cnt} indicator and therefore does not establish continuous 24-hour availability. Sixth, our outcomes are forecasting metrics and ranks rather than downstream inventory cost or service levels.

These limitations bound the claim: the paper demonstrates an evaluation mechanism and its heterogeneity, not a universal law of retail forecasting.

\section{Conclusion}

Retail forecasting evaluation can depend on a seemingly mundane preprocessing decision: whether history begins at a grounded lifecycle/observation start or at the first positive sale. Across Visuelle 2.0 and FreshRetailNet-50K, we find that this decision can alter local forecast errors, model ordering, and the set of eligible early forecast origins. The effects are not uniformly harmful and are strongly context-dependent. Visuelle exhibits broad model-level degradation and substantial rank reordering, while FreshRetailNet-50K shows smaller mixed-sign effects and non-monotonic temporal heterogeneity.

At the same time, both datasets show strong benchmark-level robustness, consistent with affected histories being uncommon. This contrast is the central empirical result: \emph{local rank sensitivity does not necessarily imply global benchmark invalidity}. Finally, scaled metrics whose denominators depend on the available history can introduce a second measurement channel and even reverse the apparent sign of a preprocessing effect.

History-start conventions should therefore be treated as explicit components of forecast evaluation design. Reporting their semantics, prevalence, eligibility consequences, and metric interactions can make retail forecasting benchmarks more transparent and easier to interpret.

\section*{Data and Code Availability}

Visuelle 2.0 is obtained from the official project page, which describes the dataset as publicly available and provides access through a short form \citep{visuelle_dataset_page}. We do not redistribute the raw Visuelle 2.0 files. FreshRetailNet-50K is obtained from the official Dingdong-Inc Hugging Face release \citep{freshretail_dataset_card}; the dataset card lists version 1.0 and a Creative Commons Attribution 4.0 International (CC BY 4.0) license. We use only the training split and do not redistribute raw FreshRetailNet-50K files in the code package. Users should obtain both datasets from their original providers and follow the corresponding terms and citation requests.

The reproducibility package is publicly available in the \href{https://github.com/hjy075/history-start-sensitivity}{GitHub repository}. It contains the analysis notebooks, version-controlled experiment logic, paper source, derived summary tables, and code that regenerates all paper figures; no third-party raw data are redistributed. The Visuelle workflow intentionally requires a user-supplied dataset path obtained through the official access route, while the FreshRetail workflow downloads the official Hugging Face dataset by repository identifier.

\section*{Reproducibility}

All experiments use public datasets and canonical implementations from StatsForecast 2.1.1. The analysis was executed in Google Colab with persistent data caches and checkpointed intermediate outputs. The final uncertainty analysis uses 5,000 product-cluster bootstrap replicates while targeting the same series-weighted estimand as the point estimates.

\bibliographystyle{plainnat}
\bibliography{references}
\appendix
\section{Additional Metric-Sensitivity Figures}

\begin{figure}[h]
\centering
\begin{minipage}{0.49\linewidth}
\centering
\includegraphics[width=\linewidth]{figures/visuelle_metric_normalization.png}
\caption*{(a) Visuelle 2.0}
\end{minipage}
\hfill
\begin{minipage}{0.49\linewidth}
\centering
\includegraphics[width=\linewidth]{figures/freshretail_metric_normalization.png}
\caption*{(b) FreshRetailNet-50K}
\end{minipage}
\caption{Fixed-denominator versus convention-specific MASE sensitivity. The latter changes both the forecast and the scaling denominator.}
\label{fig:metric}
\end{figure}

\section{Robustness Notes}

For Visuelle, minimum-history values 2, 3, and 4 all preserve IMAPA as the deterministic winner while producing Kendall's $\tau$ of 0.278, 0.389, and 0.389, respectively. TSB sensitivity checks at $\alpha_d=\alpha_p\in\{0.1,0.2,0.3\}$ preserve the positive direction of the trimming effect.

For FreshRetailNet-50K, minimum-history values 2, 3, and 4 preserve TSB as the deterministic winner while changing the remaining rank order. A minimum history of 7 changes the grounded-condition winner to CrostonSBA but leaves TSB first after trimming. Tail robustness checks excluding the unusual day-18 delay concentration and restricting delay to at most seven days preserve the qualitative mechanism result.

\end{document}
